\section*{第 \ref{ch:g12:sound-acoustics} 章——声音与声学}

\begin{solution}{exo:g12:sound-acoustics:1}
$\lambda = v/f$：(a) $340/55 \approx \qty{6.2}{m}$；(b)
\qty{0.34}{m}；(c) \qty{2.3}{cm}。比值 $15000/55 \approx 270$。
\end{solution}

\begin{solution}{exo:g12:sound-acoustics:2}
\omterm{def:g10:signals-and-waves:ultrasound}{次声波}：\qty{8}{Hz}。可闻声：\qty{100}{Hz} 和 \qty{15}{kHz}。
\omterm{def:g10:signals-and-waves:ultrasound}{超声波}：\qty{40}{kHz} 和 \qty{3}{MHz}。\omterm{def:g12:sound-acoustics:pitch-loudness}{音调}较低者：
\qty{100}{Hz}（\omterm{def:g10:signals-and-waves:frequency}{频率}较低）。
\end{solution}

\begin{solution}{exo:g12:sound-acoustics:3}
钢轨：$1000/5000 = \qty{0.20}{s}$；空气：$1000/340 \approx
\qty{2.9}{s}$。两次敲击相隔 \qty{2.7}{s}，先到达的是钢轨传来的\omterm{def:g10:signals-and-waves:sound}{声音}。
\end{solution}

\begin{solution}{exo:g12:sound-acoustics:4}
(a) $L = 10\log_{10}10^{7} = \qty{70}{dB}$。(b) $I = 10^{-12} \times
10^{6} = \qty{1.0e-6}{W/m^2}$。(c) \qty{70}{dB}：吸尘器；
\qty{60}{dB}：交谈。
\end{solution}

\begin{solution}{exo:g12:sound-acoustics:5}
\qtylist{392;588;784}{Hz}。其较高的\omterm{def:g12:sound-acoustics:harmonics}{谐音}（$n \ge 2$）从
\qty{24}{kHz} 开始，高于可闻上限：只能听到 \qty{12}{kHz}
的\omterm{def:g12:sound-acoustics:harmonics}{基音}。
\end{solution}

\begin{solution}{exo:g12:sound-acoustics:6}
$I = 0.50/(4\pi \times 25) \approx \qty{1.6e-3}{W/m^2}$；
$L = 10\log_{10}(\num{1.6e9}) \approx \qty{92}{dB}$。
\end{solution}

\begin{solution}{exo:g12:sound-acoustics:7}
两台：$+\qty{3}{dB} \to \qty{71}{dB}$；四台：\qty{74}{dB}。
$+\qty{10}{dB}$ 需要强度 $\times 10$，即 $10$ 把小提琴。
\end{solution}

\begin{solution}{exo:g12:sound-acoustics:8}
$r \times 4$：$I \div 16$，因此 $-\qty{12}{dB}$：\qty{83}{dB}。是的，
刚好低于 \qty{85}{dB} 的界限。
\end{solution}

\begin{solution}{exo:g12:sound-acoustics:9}
空气：$340/440 \approx \qty{0.77}{m}$；水：$1500/440 \approx
\qty{3.4}{m}$。\omterm{def:g10:signals-and-waves:frequency}{频率}由声源决定；\omterm{def:g12:mechanical-waves:wavelength}{波长}
$\lambda = v/f$ 随介质而变。
\end{solution}

\begin{solution}{exo:g12:sound-acoustics:10}
$f_1 = \qty{130}{Hz}$（峰间距）。相同之处：各条谱线位置
$nf_1$（\omterm{def:g12:sound-acoustics:pitch-loudness}{音调}相同）以及总强度；不同之处：谱线高度
——正是这种组合形成了\omterm{def:g12:sound-acoustics:timbre}{音色}。
\end{solution}

\begin{solution}{exo:g12:sound-acoustics:11}
$v = 2Lf_1 = 2 \times 0.650 \times 110 = \qty{143}{m/s}$。八度：
$L' = v/(2 \times 220) = L/2 = \qty{32.5}{cm}$；五度：
$L' = 143/330 \approx \qty{43.3}{cm} = 2L/3$。
\end{solution}

\begin{solution}{exo:g12:sound-acoustics:12}
$\lambda = 1540/\num{5.0e6} \approx \qty{0.31}{mm}$。深度 $= vt/2$：
$1540 \times \num{26e-6}/2 \approx \qty{2.0}{cm}$ 和
$1540 \times \num{58e-6}/2 \approx \qty{4.5}{cm}$；厚度
$\approx \qty{2.5}{cm}$。
\end{solution}

\begin{solution}{exo:g12:sound-acoustics:13}
$-\qty{25}{dB}$ 意味着 $I \div 10^{2.5} \approx 316$，因此
$r \times \sqrt{316} \approx 17.8$：$r \approx 3.0 \times 17.8
\approx \qty{53}{m}$。大厅的实际情况有所不同：地面与墙壁的反射
会增加\omterm{def:g12:sound-acoustics:intensity}{声强}，而空气和听众会吸收声能（且
舞台并非各向同性的点声源）。
\end{solution}

\begin{solution}{exo:g12:sound-acoustics:14}
$1/10^{-12} = 10^{12}$。鼓膜面积 $S = \qty{5e-5}{m^2}$：听阈时的接收
\omterm{def:g11:work-of-force:power}{功率}为 $\qty{5e-17}{W}$，痛阈时为 $\qty{5e-5}{W}$——耳朵能探测到
几十阿瓦的\omterm{def:g11:work-of-force:power}{功率}，并且在\omterm{def:g12:sound-acoustics:intensity}{声强}达到其万亿倍时仍能工作：
这是一种动态范围跨越十二个\omterm{def:g10:orders-of-magnitude:oom}{数量级}的探测器。
\end{solution}

\begin{solution}{exo:g12:sound-acoustics:15}
$2$：\qty{83}{dB}；$10$：\qty{90}{dB}；$100$：\qty{100}{dB}。每
增加 \qty{+10}{dB}，扬声器的成本就要增加十倍：\omterm{def:g12:sound-acoustics:pitch-loudness}{响度}竞赛的代价
呈\omterm{def:g10:orders-of-magnitude:scinot}{指数}增长（而且在分出胜负之前，早已危险）。
\end{solution}

\begin{solution}{pb:g12:sound-acoustics:1}
\textbf{1.} $v = 2Lf_1 = 2 \times 0.650 \times 110 = \qty{143}{m/s}$。

\textbf{2.} 弦上 $\lambda_1 = 2L = \qty{1.30}{m}$；在空气中
$\lambda = 340/110 \approx \qty{3.1}{m}$。$f$ 相同，但
$\lambda = v/f$，且两者的波速不同。

\textbf{3.} $L' = v/(2 \times 220) = L/2 = \qty{32.5}{cm}$：
第十二品位于中点，距琴桥 $L/2$ 处。

\textbf{4.} $L' = 143/(2 \times 165) \approx \qty{43.3}{cm} = 2L/3$。

\textbf{5.} $\log(220/110)/\log 2^{1/12} = 12$ 个半音——
十二品即为一个八度。第一品：$L_1 = L/2^{1/12} = 65.0/1.059
\approx \qty{61.4}{cm}$。

\textbf{6.} \qtylist{110;220;330;440;550}{Hz}。

\textbf{7.} $\qty{440}{Hz} = 4f_1$：第四次谐波，
高两个八度；$\qty{330}{Hz} = 3f_1$：第三次谐波，
高一个八度加一个五度。

\textbf{8.} 拨弦位置决定振动模式的组合：靠近琴桥时
高次\omterm{def:g12:sound-acoustics:harmonics}{谐音}较强（\omterm{def:g10:light-spectra:white}{光谱}明亮），在指板上方\omterm{def:g10:signals-and-waves:oscilloscope}{时
基}音占主导（\omterm{def:g12:sound-acoustics:timbre}{音色}醇厚）。$f_1$ 相同，故\omterm{def:g12:sound-acoustics:pitch-loudness}{音调}相同——
\omterm{def:g10:light-spectra:white}{光谱}不同，\omterm{def:g12:sound-acoustics:timbre}{音色}不同。

\textbf{9.} 只有中点处为\omterm{def:g12:sound-acoustics:standing-wave}{波节}的\omterm{def:g12:sound-acoustics:harmonics}{谐音}能够保留：即
偶次\omterm{def:g12:sound-acoustics:harmonics}{谐音}，\qtylist{220;440;660}{Hz}，\dots \omterm{def:g12:sound-acoustics:pitch-loudness}{音调}
跃升一个八度，至 \qty{220}{Hz}。

\textbf{10.} $nf_1 \le \qty{20}{kHz}$：$n \le 20000/110 = 181.8$，
故有 $181$ 个\omterm{def:g12:sound-acoustics:harmonics}{谐音}。

\textbf{11.} $I = P/(4\pi r^2) = \num{6.0e-3}/(4\pi \times 9.0)
\approx \qty{5.3e-5}{W/m^2}$。

\textbf{12.} $L = 10\log_{10}(\num{5.3e7}) \approx \qty{77}{dB}$。

\textbf{13.} $I \times 4$：$+\qty{6}{dB} \to \qty{83}{dB}$。\omterm{def:g12:sound-acoustics:level}{声强级}
是对数，\omterm{def:g12:sound-acoustics:intensity}{声强}可以相加，而\omterm{def:g12:sound-acoustics:level}{声强级}不能。

\textbf{14.} $+10\log_{10}12 \approx +\qty{10.8}{dB} \to
\qty{88}{dB}$。

\textbf{15.} 目标\omterm{def:g12:sound-acoustics:intensity}{声强} $I = I_0 \times 10^{8.5} \approx
\qty{3.2e-4}{W/m^2}$，取 $12P = \qty{7.2e-2}{W}$：
$r = \sqrt{12P/(4\pi I)} \approx \sqrt{18.1} \approx \qty{4.3}{m}$。

\textbf{16.} \qty{88}{dB} 比八小时限值高出 $+\qty{3}{dB}$，
安全时间减半一次——即四小时。此场音乐会符合要求。

\textbf{17.} $\lambda = 5000/\num{2.0e6} = \qty{2.5}{mm}$。波
只有缺陷宽度至少约为一个\omterm{def:g12:mechanical-waves:wavelength}{波长}时，波才会被反射：可听声在木材中
的 $\lambda$ 为数十厘米，会
绕过任何小缺陷发生衍射。

\textbf{18.} $t = 2d/v = 2 \times 0.045/5000 = \qty{18}{\micro s}$。

\textbf{19.} $d = vt/2 = 5000 \times \num{7.2e-6}/2 = \qty{18}{mm}$。

\textbf{20.} 小结：弦上 $v = \qty{143}{m/s}$；品位在
$\qty{32.5}{cm}$（第十二品）和 $\qty{61.4}{cm}$（第一品）处；合奏
的安全距离超过 $\approx \qty{4.3}{m}$；缺陷深 \qty{18}{mm}——四
个数字，贯穿一章。
\end{solution}
